\documentclass[10pt]{article}

\usepackage[a4paper,margin=16mm,headheight=14pt]{geometry}
\usepackage[utf8]{inputenc}
\usepackage[T1]{fontenc}
\usepackage{lmodern}
\usepackage{microtype}
\usepackage{amsmath,amssymb,amsfonts,mathtools,bm}
\usepackage{booktabs,longtable,tabularx,array,multirow}
\usepackage{float}
\usepackage{enumitem}
\usepackage{xcolor}
\usepackage[round,authoryear]{natbib}
\usepackage[
  colorlinks=true,
  linkcolor=blue!45!black,
  citecolor=blue!45!black,
  urlcolor=blue!45!black
]{hyperref}
\usepackage[capitalise,nameinlink,noabbrev]{cleveref}
\usepackage{fancyhdr}

\definecolor{maincolor}{HTML}{17365D}
\definecolor{accentcolor}{HTML}{2B6F9C}
\setlength{\parindent}{0pt}
\setlength{\parskip}{0.55\baselineskip}
\setlist{nosep,leftmargin=*}
\setcounter{secnumdepth}{2}
\setcounter{tocdepth}{2}
\renewcommand{\arraystretch}{1.12}

\pagestyle{fancy}
\fancyhf{}
\fancyhead[L]{\small Retrieval-based forecast adaptation}
\fancyhead[R]{\small Gaspard Berthelier}
\fancyfoot[C]{\thepage}
\renewcommand{\headrulewidth}{0.35pt}

\newcommand{\reals}{\mathbb{R}}
\newcommand{\code}[1]{\texttt{\detokenize{#1}}}
\newcommand{\Tzero}{T_0}
\newcommand{\Tone}{T_1}
\newcommand{\Ttwo}{T_2}
\newcommand{\Tpool}{T_{12}}
\newcommand{\Tthree}{T_3}
\newcommand{\avgy}{\overline{Y}}
\newcommand{\sd}{\operatorname{sd}}
\newcommand{\MSE}{\operatorname{MSE}}
\newcommand{\nMSE}{\operatorname{nMSE}}
\newcommand{\TopK}{\operatorname{TopK}}

\title{\vspace{-1.5em}\textbf{Retrieval-Based Adaptation of \\Time Series Foundation Models}\\[2mm]\large Experiment Guideline\vspace{-0.4em}}
\author{Gaspard Berthelier}
\date{7 August 2026}

\begin{document}
\maketitle

\tableofcontents

\section{Time Series Forecasting}
\label{sec:forecasting}

Let $\mathcal U=\{0,\ldots,U-1\}$ be a set of $U$ users or channels and
$t\in\{0,\ldots,T-1\}$ the time index for the $T$ measured dates. For channel
$u\in\mathcal U$, the univariate series is
$z^{(u)}_0,\ldots,z^{(u)}_{T-1}$. We consider the setting of a fixed lookback length $L\geq1$ and forecast
horizon $H\geq1$. Given a valid \emph{query date} $s$, the lookback and future horizon are
\begin{align}
 X_{s,u}
 &= z^{(u)}_{(s-L,s]}
  = (z^{(u)}_{s-L+1},\ldots,z^{(u)}_s)\in\reals^L,\\
 Y_{s,u}
 &= z^{(u)}_{(s,s+H]}
  = (z^{(u)}_{s+1},\ldots,z^{(u)}_{s+H})\in\reals^H.
\end{align}
Thus the lookback length $L$ represents the number of observed values at a given time and the horizon $H$ is the number of target values to predict.

A forecasting model is a parameterized map
\[
 F_\theta:\reals^L\longrightarrow\reals^H,
 \qquad X_{s,u}\longmapsto \widehat Y_{s,u}.
\]
Here $\theta$ denotes the model parameters.

Certain families of models allow the use of optional covariates. Let $C_{s,u}$ denote a set of covariate windows of size $L+H$ (horizon included).
Whe using covariates, the prediction becomes $\widehat Y^{\rm cov}_{s,u}=F_\theta(X_{s,u};C_{s,u})$.

For any $x=(x_1,\ldots,x_n)\in\reals^n$, define
\[
 \mu_x=\frac1n\sum_{k=1}^n x_k,
 \qquad
 \sigma_x=\sd(x)=\left(\frac1n\sum_{k=1}^n(x_k-\mu_x)^2\right)^{1/2},
 \qquad \varepsilon=10^{-8}.
\]

For one sample, consisting of windows $(X,Y,\widehat Y)$, the principal losses are
\[
 \MSE(\widehat Y,Y)=\frac1H\sum_{h=1}^H(\widehat Y_h-Y_h)^2,
 \qquad
 \nMSE(\widehat Y,Y\mid X)
 =\frac1H\sum_{h=1}^H
 \left(\frac{\widehat Y_h-Y_h}{\max(\sigma_X,\varepsilon)}\right)^2.
\]

\section{Foundation Models}
\label{sec:foundation-models}

A time series Foundation Model (FM) is pretrained over many tasks and datasets, and directly reused on
unseen series without updating its parameters. Indeed, their extensive training allows them to perform In-Context Learning (ICL),
which means the model has learned to detect patterns within the input, as if it were a dataset on which it was fine-tuned.
A few examples of FMs are provided below.

\textbf{Chronos-2.} Chronos-2 is a 120M-parameter encoder-only foundation model
whose group attention supports univariate, multivariate, and covariate-aware
zero-shot forecasting \citep{chronos2}. In this project, we notably use context mode
\code{future_included}: each covariate contributes its lookback as well as its observed horizon. The active
code name is \code{chronos2}.

\textbf{TabPFN-TS.} TabPFN-TS casts forecasting as tabular regression with time,
seasonal, frequency, and optional covariate features, then performs ICL with a
pretrained TabPFN regressor \citep{tabpfnts}. The project wrapper uses the
\code{tabpfn-v2.5-regressor-v2.5_default.ckpt} checkpoint, automatically adds
periods 24 and 168 when $L$ permits, and exposes the code name \code{tabpfnts}.

\textbf{TS-ICL.} TS-ICL is a 27M-parameter time-indexed encoder--regressor
Transformer for forecasting and imputation on regular or irregular timestamps
\citep{tsicl}. It is part of the literature scope, but \code{ts_icl} is
deliberately rejected by the current launcher because no wrapper is registered.

\textbf{Chronos-Bolt.} The \code{chronos-bolt} wrapper supplies the frozen
205M-parameter Chronos-Bolt-base backbone for the Cross-RAG comparison. It is
univariate and ignores retrieved covariates; neighbour targets and forecasts
remain available to the external adapter \citep{chronos,chronoscode}.

\begin{table}[H]
\centering
\small
\caption{Model names present in the Python registry or publication launchers.
Trainable means parameters updated by this project, not total pretrained size.}
\label{tab:model-registry}
\begin{tabularx}{\linewidth}{@{}lXrr@{}}
\toprule
Code name & Role and current status & Total parameters & Trainable here \\
\midrule
\code{chronos2} & Primary screen/full backbone; frozen & 120M & 0 \\
\code{tabpfnts} & Ultra backbone; frozen TabPFN-v2.5 regressor & $\approx11$M & 0 \\
\code{chronos-bolt} & Fixed Cross-RAG backbone; frozen & 205M & 0 \\
\code{persistence} & Python-registry diagnostic & 0 & 0 \\
\code{linear} & Python-registry diagnostic, univariate & $H(L+1)$ & $H(L+1)$ \\
\code{patchtst} & Registry diagnostic; not launched for publication
& 466,843 / 663,091 / 1,762,603 & same \\
\code{ts_icl} & Reserved shell name; not registered & 27M in the paper & not run \\
\bottomrule
\end{tabularx}
\end{table}

The three PatchTST numbers in \cref{tab:model-registry} use the code defaults
at $168{:}24$, $336{:}48$, and $504{:}168$, respectively. Publication
experiments accept only the three frozen backbones listed first.

\section{From RAG to retrieval-augmented forecasting}
\label{sec:history}

Retrieval-augmented generation separates a parametric model from an explicit,
replaceable non-parametric memory. Classical RAG retrieves documents from a
dense index and conditions a pretrained generator on them
\citep{lewis2020rag}. In $k$NN-MT, a frozen neural machine-translation model
caches decoder representations $k_i$ with next-token values $v_i$; at decoding
state $q_t$, it interpolates the model and neighbour distributions
\citep{khandelwal2021knnmt}. Let $p_\theta(\cdot\mid q_t)$ be the neural
next-token distribution for candidate token $y_t$, let $\mathcal N_K(q_t)$
contain the $K\geq1$ cached keys
nearest to $q_t$ under distance $d$, and choose interpolation weight
$\lambda\in[0,1]$ and temperature $\tau>0$. Then
\[
 p(y_t\mid q_t)
 =(1-\lambda)p_\theta(y_t\mid q_t)
 +\lambda\frac{\sum_{i\in\mathcal N_K(q_t)}
 \mathbf1[v_i=y_t]\exp(-d(q_t,k_i)/\tau)}
 {\sum_{i\in\mathcal N_K(q_t)}\exp(-d(q_t,k_i)/\tau)}.
\]
This is the direct ancestor of our question: can a frozen pretrained model be
adapted cheaply by retrieving domain examples and mixing their outcomes?

Time-series work has followed several paths: retrieved sequences can become a
text prompt (TimeRAG), an appended context (RAF), an input to a model trained on
the target dataset (RAFT), or values fused by pretrained attention modules
(TS-RAG and Cross-RAG). Our method is narrower: the datastore is target-panel
history, the backbone is frozen, and simple validated forecasts or gates are
tested before a neural adapter.

\small
\begin{longtable}{@{}p{.12\linewidth}p{.13\linewidth}p{.25\linewidth}p{.42\linewidth}@{}}
\caption{Bibliography map for pretrained retrieval and time-series forecasting.
This is a method map, not a cross-paper accuracy ranking.}
\label{tab:bibliography-map}\\
\toprule
Method & Date / venue & Retrieval object & Adaptation mechanism \\
\midrule
\endfirsthead
\toprule
Method & Date / venue & Retrieval object & Adaptation mechanism \\
\midrule
\endhead
RAG \citep{lewis2020rag}
& 2020, NeurIPS
& Dense-indexed passages
& Fine-tuned seq2seq generator combines parametric and non-parametric memory. \\
$k$NN-MT \citep{khandelwal2021knnmt}
& 2021, ICLR
& Decoder state / next-token pairs
& No retraining; interpolates neural and neighbour token distributions. \\
RAF \citep{tire2024raf}
& 2024 preprint; v2 2026
& Related time-series contexts and futures
& Normalises, aligns, and prepends retrieved examples to frozen TSFMs. \\
TimeRAG \citep{yang2024timerag}
& 2024 preprint; ICASSP 2025
& DTW-similar historical sequences
& Renders references and query as an LLM forecasting prompt. \\
TS-RAG \citep{ning2025tsrag}
& 2025, NeurIPS
& Encoder-space context--future pairs
& Pretrains projectors and an Adaptive Retrieval Mixer; no target-task tuning. \\
RAFT \citep{han2025raft}
& 2025, ICML
& Correlated training inputs and following patches
& Trains a shallow multiresolution forecasting MLP per target dataset. \\
Cross-RAG \citep{lee2026crossrag}
& 2026, KDD MILETS
& Min--max lookback/future pairs, cosine search
& Pretrained query--retrieval cross-attention with query-dependent and
query-independent paths. \\
This project
& 2026
& Raw or transformed historical lookback/future pairs
& Closed-form baselines, CatBoost gates, and TS-IFA around a frozen backbone. \\
\bottomrule
\end{longtable}
\normalsize

\section{Neighbour retrieval}
\label{sec:retrieval}

\subsection{Datastore and exact search}

Fix a named retrieval transform $\rho$ and distance metric $\delta$. Their
maps are
\[
 \phi=\phi_\rho:\reals^L\to\reals^{d_\rho},
 \qquad
 d=d_\delta:\reals^{d_\rho}\times\reals^{d_\rho}\to\reals_{\geq0},
\]
with every implemented $\phi_\rho$ and $d_\delta$ defined explicitly in
\cref{tab:retrieval-spaces}. Here $d_\rho=L$ except for the encoder space;
$g_\theta:\reals^L\to\reals^{d_{\rm enc}}$ denotes the frozen backbone
representation after the configured flattening or pooling.

For query $(s,u)$, let
$\mathcal R_s=\{r_{s,0}<\cdots<r_{s,n_s^{\rm store}-1}\}$ be the ordered finite
set of $n_s^{\rm store}=|\mathcal R_s|$ store dates in
$\{L-1,\ldots,T-H-1\}$ that pass the fixed/online, period, stride, and cap
rules in \cref{sec:practical}. The eligible store keys and the resulting
labelled datastore are
\[
 \mathcal I_s=\mathcal R_s\times\mathcal U,
 \qquad
 \mathcal D_s
 =\{(\phi(X_{r,v}),X_{r,v},Y_{r,v},r,v):(r,v)\in\mathcal I_s\}.
\]
For $K\in\{1,\ldots,|\mathcal I_s|\}$, let $\TopK^{\min}$ return the ordered
$K$ minimizers. Exact search returns the following store keys:
\[
 \mathcal N_K(s,u)
 =\TopK_{(r,v)\in\mathcal I_s}^{\min}
 d\!\left(\phi(X_{s,u}),\phi(X_{r,v})\right).
\]
The datastore is flattened in \emph{user-major, date-minor} order. For
$i\in\{0,\ldots,Un_s^{\rm store}-1\}$, a flat index decodes as
\[
 v=\left\lfloor\frac{i}{n_s^{\rm store}}\right\rfloor,
 \qquad m=i\bmod n_s^{\rm store},
 \qquad r=r_{s,m}.
\]
Downstream query payloads use the opposite, date-major ordering
$(s_0,u_0),\ldots,(s_0,u_{U-1}),(s_1,u_0),\ldots$ so that whole dates form
contiguous blocks for chronological splitting; $s_0<s_1<\cdots$ are the saved
query dates and $u_k=k$ is the internal channel index.

In \cref{tab:retrieval-spaces}, $X\in\reals^L$ is a generic lookback; in its
distance rows, $a,b\in\reals^{d_\rho}$ are two transformed lookbacks and
$x_{\min}=\min_{1\leq k\leq L}X_k$,
$x_{\max}=\max_{1\leq k\leq L}X_k$,
$a^\circ=a-\mu_a\mathbf1_{d_\rho}$,
$b^\circ=b-\mu_b\mathbf1_{d_\rho}$ are their centred forms.

\begin{table}[H]
\centering
\small
\caption{Implemented retrieval spaces and distances.}
\label{tab:retrieval-spaces}
\begin{tabularx}{\linewidth}{@{}llX@{}}
\toprule
Space & $\phi(X)$ & Status \\
\midrule
Raw & $X$ & Primary screen \\
Instance & $(X-\mu_X\mathbf1_L)/(\sigma_X+\varepsilon)$ & Primary screen \\
Min--max & $(X-x_{\min}\mathbf1_L)/\max(x_{\max}-x_{\min},\varepsilon)$ & Cross-RAG profile \\
Fourier & $|\operatorname{FFT}((X-\mu_X\mathbf1_L)/(\sigma_X+\varepsilon))|$ & Optional ablation \\
Encoder & flattened or pooled frozen representation $g_\theta(X)$ & Implemented override \\
\midrule
Euclidean & $\|a-b\|_2$ & Primary metric \\
Cosine & $1-a^\top b/[\max(\|a\|_2,\varepsilon)\max(\|b\|_2,\varepsilon)]$ & Cross-RAG metric \\
Pearson & $1-(a^\circ)^\top b^\circ/[\max(\|a^\circ\|_2,\varepsilon)\max(\|b^\circ\|_2,\varepsilon)]$ & Implemented override \\
\bottomrule
\end{tabularx}
\end{table}

\subsection{Saved information and query-scale transfer}

For neighbour rank $j\in\{1,\ldots,K\}$ with
$(r_j,v_j)\in\mathcal N_K(s,u)$, abbreviate
$X_j=X_{r_j,v_j}$ and $Y_j=Y_{r_j,v_j}$, and let $d_j$ be its retrieval
distance and $N_j=F_\theta(X_j)$ its frozen forecast. Also abbreviate the query
as $X=X_{s,u}$ and $Y=Y_{s,u}$. Extraction saves the quantities in
\cref{tab:retrieval-payload}; it stores $E_j=Y_j-N_j$ rather than duplicating
$N_j$, which is reconstructed downstream.

\begin{table}[H]
\centering
\small
\caption{Current extraction payload. Optional fields are generated only when
their switches are enabled.}
\label{tab:retrieval-payload}
\begin{tabularx}{\linewidth}{@{}lX@{}}
\toprule
Object & Saved content \\
\midrule
Query & $X,Y,V=F_\theta(X),C=F_\theta(X;\{X_j,Y_j\}_{j=1}^K)$ \\
Neighbours & $X_j,Y_j,E_j=Y_j-N_j,d_j,r_j,v_j$ for $j=1,\ldots,K$ \\
Indices & query date/user, neighbour date/user, eligible store-date/window counts \\
Summaries & query and neighbour level/scale; forecast--context and residual losses \\
Optional & context-conditioned neighbour residuals \\
\bottomrule
\end{tabularx}
\end{table}

Retrieved values are expressed on the query scale before fitting. Write
$q=(s,u)$ for the current query and let
$(\mu_q,\sigma_q)=(\mu_X,\sigma_X)$ and
$(\mu_j,\sigma_j)=(\mu_{X_j},\sigma_{X_j})$. For a neighbour-scale horizon
vector $Z\in\reals^H$ or residual vector $E\in\reals^H$, define
\begin{align}
 \mathcal T_{j\to q}(Z)
 &=\frac{Z-\mu_j\mathbf1_H}{\max(\sigma_j,\varepsilon)}\sigma_q
   +\mu_q\mathbf1_H,\\
 \mathcal T^{\rm res}_{j\to q}(E)
 &=\frac{E}{\max(\sigma_j,\varepsilon)}\sigma_q.
\end{align}
All later symbols $Y_j,N_j,E_j$ denote these query-scale values.

Distances are converted to weights independently for every query. Let
$\bm d=(d_1,\ldots,d_K)$ be the ordered retrieved-distance vector and let
$j,\ell\in\{1,\ldots,K\}$:
\[
 \widetilde d_j
 =\frac{d_j-\min_\ell d_\ell}{\max(\sigma_{\bm d},\varepsilon)},
 \qquad
 w_j=\frac{e^{-\widetilde d_j}}{\sum_{\ell=1}^K e^{-\widetilde d_\ell}},
 \qquad
 \avgy_h=\sum_{j=1}^K w_jY_{j,h}.
\]
Thus $w_j\geq0$, $\sum_{j=1}^K w_j=1$, and $\avgy\in\reals^H$; for $K=1$,
$w_1=1$.

\section{Retrieval baselines}
\label{sec:baselines}

For extracted window $i$ with query $(s_i,u_i)$, define the target
$G_i=Y_{s_i,u_i}\in\reals^H$, vanilla forecast $V_i\in\reals^H$,
conditioned forecast $C_i\in\reals^H$, query-scale neighbour future and
forecast $(Y_{i,j},N_{i,j})\in\reals^H\times\reals^H$ for
$j\in\{1,\ldots,K\}$, and
\[
 V_i=F_\theta(X_{s_i,u_i}),\qquad
 C_i=F_\theta(X_{s_i,u_i};\{X_{i,j},Y_{i,j}\}_{j=1}^K),\qquad
 \avgy_i=\sum_{j=1}^K w_{i,j}Y_{i,j}.
\]
Each design is $(V_i,Z_{i,1},\ldots,Z_{i,p})$.

\begin{table}[H]
\centering
\small
\caption{Baseline designs; $Z_i=(Z_{i,1},\ldots,Z_{i,p})$.}
\label{tab:baseline-designs}
\begin{tabular}{@{}lll@{}}
\toprule
Stem & $Z_i$ & $p$ \\
\midrule
\code{cov} & $(C_i)$ & 1 \\
\code{avgy} & $(\avgy_i)$ & 1 \\
\code{y} & $(Y_{i,1},\ldots,Y_{i,K})$ & $K$ \\
\code{cov_y} & $(C_i,Y_{i,1},\ldots,Y_{i,K})$ & $1+K$ \\
\code{cov_avgy} & $(C_i,\avgy_i)$ & 2 \\
\code{residual} & $(Y_{i,1},\ldots,Y_{i,K},N_{i,1},\ldots,N_{i,K})$ & $2K$ \\
\code{full} & $(C_i,Y_{i,1},\ldots,Y_{i,K},N_{i,1},\ldots,N_{i,K})$ & $1+2K$ \\
\bottomrule
\end{tabular}
\end{table}

Let $\tau\in\{S,H\}$ denote shared or horizon-wise coefficients. For any
coefficient vector $c$, set $c_h^S=c$ for every $h$, whereas
$c_1^H,\ldots,c_H^H$ are fitted independently. Define the simplex
$\mathcal S_p=\{\lambda\in\reals^{p+1}:\lambda_m\geq0,
\sum_{m=0}^{p}\lambda_m=1\}$. The three families are
\begin{align}
 \widehat G^{R,\tau}_{i,h}
 &=V_{i,h}+a_h^\tau V_{i,h}
   +\sum_{m=1}^{p}b_{h,m}^\tau Z_{i,h,m},
 &&(a_h^\tau,b_{h,1}^\tau,\ldots,b_{h,p}^\tau)\in\reals^{p+1},
 \label{eq:ridge-family}\\
 \widehat G^{C,\tau}_{i,h}
 &=\lambda_{h,0}^\tau V_{i,h}
   +\sum_{m=1}^{p}\lambda_{h,m}^\tau Z_{i,h,m},
 &&\lambda_h^\tau\in\mathcal S_p,
 \label{eq:convex-family}\\
 \widehat G^{D,\tau}_{i,h}
 &=V_{i,h}+\sum_{m=1}^{p}b_{h,m}^\tau
   (Z_{i,h,m}-V_{i,h}),
 &&b_h^\tau\in\reals^p.
 \label{eq:delta-family}
\end{align}
Vanilla is recovered by $(a,b)=0$, $\lambda=(1,0,\ldots,0)$, or $b=0$.
Ridge contains every convex forecast through
$a=\lambda_0-1$ and $b_m=\lambda_m$. For \code{cov}, the formulas are
$V+aV+bC$, $\lambda V+(1-\lambda)C$, and $V+b(C-V)$.

Let $\mathcal Q_A$ be the windows in split $A\in\{\Tone,\Ttwo,\Tpool\}$,
where $\Tpool=\Tone\cup\Ttwo$. For $\tau=S$, set
$\mathcal O_A^S=\mathcal Q_A\times\{1,\ldots,H\}$; for the $h$th
horizon fit, set $\mathcal O_A^{H,h}=\mathcal Q_A\times\{h\}$. Ridge and
delta-ridge use correction target $t_{i,h}=G_{i,h}-V_{i,h}$ and design
\[
 d^R_{i,h}=(V_{i,h},Z_{i,h,1},\ldots,Z_{i,h,p})\in\reals^{p+1},
 \qquad
 d^D_{i,h}=(Z_{i,h,1}-V_{i,h},\ldots,Z_{i,h,p}-V_{i,h})\in\reals^p.
\]
For either $d\in\reals^q$ and an observation set $\mathcal O$, define
\[
 r_k=\max\!\left(
  \sqrt{\frac1{|\mathcal O|}\sum_{(i,h)\in\mathcal O}d_{i,h,k}^2},10^{-12}
 \right),
 \qquad \widetilde d_{i,h,k}=d_{i,h,k}/r_k,
\]
and, for $\alpha\geq0$,
\[
 \widehat\gamma_{\alpha,\mathcal O}
 =\arg\min_{\gamma\in\reals^{q}}
 \frac1{|\mathcal O|}\sum_{(i,h)\in\mathcal O}
 (t_{i,h}-\widetilde d_{i,h}^{\top}\gamma)^2
 +\alpha\|\gamma\|_2^2,
 \qquad
 \widehat\beta_{\alpha,\mathcal O,k}
 =\widehat\gamma_{\alpha,\mathcal O,k}/r_k.
\]
The grid is
\[
 \mathcal A=\{0,10^{-6},10^{-5},10^{-4},10^{-3},10^{-2},10^{-1},1,10\}.
\]
For each $\alpha\in\mathcal A$, coefficients fitted on $\Tone$ are scored by
nMSE on $\Ttwo$; the minimizer (smallest $\alpha$ on ties) is refitted on
$\Tpool$. Convex coefficients are fitted directly on $\Tpool$:
\[
 \widehat\lambda_{\mathcal O}
 =\arg\min_{\lambda\in\mathcal S_p}
 \frac1{|\mathcal O|}\sum_{(i,h)\in\mathcal O}
 \left(G_{i,h}-\lambda_0V_{i,h}
 -\sum_{m=1}^p\lambda_mZ_{i,h,m}\right)^2.
\]
Use $\mathcal O=\mathcal O_A^S$ for a shared fit and
$\mathcal O=\mathcal O_A^{H,h}$ independently for each $h$; no intercept or
centring is added. Chunked sufficient statistics give the exact objectives.

\small
\begin{longtable}{@{}p{.30\linewidth}p{.46\linewidth}p{.16\linewidth}@{}}
\caption{Baseline names; substitute a stem from \cref{tab:baseline-designs}.}
\label{tab:baselines}\\
\toprule
Artifact & Definition & Fitted scalars \\
\midrule
\endfirsthead
\toprule
Artifact & Definition & Fitted scalars \\
\midrule
\endhead
\code{vanilla} & $V_{i,h}$ & 0 \\
\code{cov_forecast} & $C_{i,h}$ & 0 \\
\code{avgy} & $\avgy_{i,h}$ & 0 \\
\code{y_mean} & $K^{-1}\sum_jY_{i,j,h}$ & 0 \\
\midrule
\code{<stem>_ridge_*} & \cref{eq:ridge-family}, $*=\code{shared}/\code{horizon}$ & $(p+1)$ / $H(p+1)$ \\
\code{<stem>_convex_*} & \cref{eq:convex-family}, $*=\code{shared}/\code{horizon}$ & $(p+1)$ / $H(p+1)$ \\
\code{<stem>_delta_ridge_*} & \cref{eq:delta-family}, $*=\code{shared}/\code{horizon}$ & $p$ / $Hp$ \\
\bottomrule
\end{longtable}
\normalsize

\section{Gates and features}
\label{sec:gates}

A gate compares $V_i$ with $A_i=C_i$ (\code{cov}) or
$A_i=\avgy_i$ (\code{avgy}). Define the pointwise and shared candidate
advantages
\[
 \Delta_{i,h}=(G_{i,h}-V_{i,h})^2-(G_{i,h}-A_{i,h})^2,
 \qquad \overline\Delta_i=H^{-1}\sum_{h=1}^H\Delta_{i,h}.
\]

For $q=(q_1,\ldots,q_K)\in\reals^K$, define neighbour-weighted moments
\[
 \mu_{w_i}(q)=\sum_{j=1}^K w_{i,j}q_j,\qquad
 \sigma_{w_i}(q)=
 \left[\sum_{j=1}^K w_{i,j}(q_j-\mu_{w_i}(q))^2\right]^{1/2}.
\]
Let $\mu_h$ and $\sigma_h$ denote population moments over
$h\in\{1,\ldots,H\}$. For neighbour $j$, set
\begin{align*}
 r^V_{i,j}&=\mu_h(Y_{i,j,h}-V_{i,h}),
 &s^V_{i,j}&=\sigma_h(Y_{i,j,h}-V_{i,h}),\\
 r^N_{i,j}&=\mu_h(Y_{i,j,h}-N_{i,j,h}),
 &s^N_{i,j}&=\sigma_h(Y_{i,j,h}-N_{i,j,h}).
\end{align*}
Let $(r_{i,j},v_{i,j})$ be neighbour $j$'s date and channel, and define
\[
 \bm\ell_i=(\mu_{X_{i,1}},\ldots,\mu_{X_{i,K}}),\quad
 \bm\rho_i=(\sigma_{X_{i,1}},\ldots,\sigma_{X_{i,K}}),\quad
 \bm a_i=(s_i-r_{i,1},\ldots,s_i-r_{i,K}),
\]
\[
 \bm w_i=(w_{i,1},\ldots,w_{i,K}),\qquad
 \bm d_i=(d_{i,1},\ldots,d_{i,K}).
\]
The exact static vector $f_i^0\in\reals^{18}$ is
\begin{align*}
f_i^0=(&
 \mu_{w_i}(\bm r_i^V),\sigma_{w_i}(\bm r_i^V),\mu_{w_i}(\bm s_i^V),
 \mu_{w_i}(\bm r_i^N),\sigma_{w_i}(\bm r_i^N),\mu_{w_i}(\bm s_i^N),\\
&\mu_{X_{s_i,u_i}},\sigma_{X_{s_i,u_i}},
 \mu_{\bm\ell_i},\sigma_{\bm\ell_i},\mu_{\bm\rho_i},\sigma_{\bm\rho_i},
 K^{-1}\textstyle\sum_{j=1}^K\mathbf1[v_{i,j}=u_i],\\
&\mu_{\bm a_i},\sigma_{\bm a_i},\sigma_{\bm w_i},
 \max_j w_{i,j},\mu_{\bm d_i}),
\end{align*}
where $\bm r_i^V=(r^V_{i,j})_{j=1}^K$ and analogously for
$\bm s_i^V,\bm r_i^N,\bm s_i^N$. The shared and horizon feature vectors are
\begin{align}
 f_i^S={}&(\mu_h(A_{i,h}-V_{i,h}),\sigma_h(A_{i,h}-V_{i,h}),f_i^0)
 \in\reals^{20},\label{eq:shared-gate-features}\\
 f_{i,h}^H={}&(A_{i,h}-V_{i,h},
 \mu_{w_i}((Y_{i,j,h}-V_{i,h})_j),
 \sigma_{w_i}((Y_{i,j,h}-V_{i,h})_j),\notag\\
 &\mu_{w_i}((Y_{i,j,h}-N_{i,j,h})_j),
 \sigma_{w_i}((Y_{i,j,h}-N_{i,j,h})_j),f_i^0)
 \in\reals^{23}.\label{eq:horizon-gate-features}
\end{align}

For shape $\tau\in\{S,H\}$, write
$z_i^S=\overline\Delta_i$, $z_{i,h}^H=\Delta_{i,h}$,
$f_i^S$ as in \cref{eq:shared-gate-features}, and $f_{i,h}^H$ as in
\cref{eq:horizon-gate-features}. A shared model is fitted once; a horizon
model is fitted independently for each $h$. CatBoost scores are
\[
 g^{\mathrm{reg},\tau}=\operatorname{CBReg}(f^\tau),\qquad
 g^{\mathrm{cls},\tau}
 =\Pr_{\operatorname{CBCls}}(z^\tau>0\mid f^\tau)-\tfrac12.
\]
The tree count is selected on $\Ttwo$ after fitting on $\Tone$, then the model
is refitted on $\Tpool$ with that count. Define
\[
 q^{\mathrm{hard},\tau}=\mathbf1[g^\tau>0],\qquad
 \widehat G^{\mathrm{hard},\tau}_{i,h}
 =(1-q_{i,h}^{\mathrm{hard},\tau})V_{i,h}
 +q_{i,h}^{\mathrm{hard},\tau}A_{i,h},
\]
with $q_{i,h}^{\cdot,S}=q_i^{\cdot,S}$. For soft output, let
$\operatorname{sigmoid}(x)=(1+e^{-x})^{-1}$ and
\[
 s^S=\sd((z_i^S)_{i\in\mathcal Q_{\Tpool}}),\qquad
 s_h^H=\sd((z_{i,h}^H)_{i\in\mathcal Q_{\Tpool}}).
\]
Writing $s^\tau$ for the matching shared or $h$th-horizon scale,
\[
 p^{\mathrm{cls},\tau}=g^{\mathrm{cls},\tau}+\tfrac12,\qquad
 p^{\mathrm{reg},\tau}=
 \begin{cases}
 \operatorname{sigmoid}(g^{\mathrm{reg},\tau}/s^\tau),&s^\tau>10^{-12},\\
 0,&s^\tau\leq10^{-12},\ g^{\mathrm{reg},\tau}<0,\\
 \tfrac12,&s^\tau\leq10^{-12},\ g^{\mathrm{reg},\tau}=0,\\
 1,&s^\tau\leq10^{-12},\ g^{\mathrm{reg},\tau}>0,
 \end{cases}
\]
\[
 \widehat G^{\mathrm{soft},\tau}_{i,h}
 =(1-p_{i,h}^\tau)V_{i,h}+p_{i,h}^\tau A_{i,h}.
\]
For shared output, $p_{i,h}^{\cdot,S}=p_i^{\cdot,S}$ for every $h$.
The no-feature Bayes scores and scored-split oracle scores are
\[
 g^{B,S}=|\mathcal Q_{\Tpool}|^{-1}\sum_{i\in\mathcal Q_{\Tpool}}
 \overline\Delta_i,\quad
 g_h^{B,H}=|\mathcal Q_{\Tpool}|^{-1}\sum_{i\in\mathcal Q_{\Tpool}}
 \Delta_{i,h},\quad
 g_i^{O,S}=\overline\Delta_i,\quad g_{i,h}^{O,H}=\Delta_{i,h},
\]
and use the hard rule above.

\begin{table}[H]
\centering
\small
\caption{Gate names for $A\in\{\code{cov},\code{avgy}\}$; $\tau=S/H$.}
\label{tab:gate-models}
\begin{tabularx}{\linewidth}{@{}lXr@{}}
\toprule
Name pattern & Definition & Fitted scalars / proxy \\
\midrule
\code{bayes_<A>_shared/horizon} & $g^{B,\tau}$, hard & $1$ / $H$ \\
\code{catboost_<A>_regressor_shared/horizon} & $g^{\mathrm{reg},\tau}$, hard & $\leq6000$ / $\leq6000H$ \\
\code{catboost_<A>_classifier_shared/horizon} & $g^{\mathrm{cls},\tau}$, hard & $\leq6000$ / $\leq6000H$ \\
\code{catboost_<A>_*_shared/horizon_soft} & same fit, soft output & 0 additional \\
\code{oracle_<A>_shared/horizon} & $g^{O,\tau}$, hard & 0 \\
\bottomrule
\end{tabularx}
\end{table}

CatBoost uses at most 300 depth-4 trees, learning rate 0.03, 50-round early
stopping, and balanced classifier classes \citep{prokhorenkova2018catboost}.
The proxy is $300(16\text{ leaf}+4\text{ split})=6000$. Screen uses only hard
shared regressors; classifier, soft-regressor, and horizon-regressor forms are
one-axis ablations.

\section{Practical protocol}
\label{sec:practical}

\subsection{Chronological regions and window eligibility}

For the $T$ dates introduced in \cref{sec:forecasting}, extraction computes
exclusive boundaries, where $\operatorname{round}$ means nearest-integer
rounding and $[a,b)=\{a,a+1,\ldots,b-1\}$ for integer dates:
\[
 b_0=\operatorname{round}(0.30T),\qquad
 b_{12}=\operatorname{round}(0.80T),\qquad b_3=T.
\]
The extraction regions are the fixed datastore region $\Tzero=[0,b_0)$,
pooled adaptation region $\Tpool=[b_0,b_{12})$, and held-out evaluation
region $\Tthree=[b_{12},T)$. For a target period
$[b,e)$, eligible query dates are
\[
 s\in
 \left[\max(L-1,b-1),\ \min(T-H-1,e-H-1)\right]
\]
at the configured query stride. Hence the first target may start at $b$ with
query index $b-1$; its lookback may cross the boundary, but its complete target
must be inside $[b,e)$. Changing $L$ therefore preserves $\Tthree$ target dates
whenever both runs remain eligible.

Each downstream learner re-splits $\Tpool$ by unique query date. If
$n_{12}\geq2$ dates are present and the validation fraction is $\rho=0.2$,
\[
 n_2=\min(n_{12}-1,\max(1,\operatorname{round}(\rho n_{12}))),
 \qquad n_1=n_{12}-n_2.
\]
The first $n_1$ and last $n_2$ whole dates form $\Tone$ and $\Ttwo$,
respectively, so $\Tpool=\Tone\cup\Ttwo$ and the two sets are disjoint; users
on the same date are never split.

\subsection{Datastore date grids}

Let $P$ be the retrieval period, $D_{\rm stride}>0$ the datastore stride, and
$A\in\{0,1\}$ the period-alignment indicator. Alignment requires
$P\mid D_{\rm stride}$. Before optional explicit history bounds, define the
admissible endpoint pair
\[
 (a_s,b_s)=
 \begin{cases}
  (L-1,b_0-H-1),&\text{fixed mode},\\
  (L-1,s-H),&\text{online mode}.
 \end{cases}
\]
The first candidate and uncapped date grid are
\[
 r_s^{A}=\begin{cases}
 a_s+((s-a_s)\bmod P),&A=1,\\
 a_s,&A=0,
 \end{cases}
 \qquad
 \mathcal G_s=\{r_s^{A}+\ell D_{\rm stride}:\ell\in\mathbb N_0,\
 r_s^{A}+\ell D_{\rm stride}\leq b_s\}.
\]
Let $\operatorname{Tail}_M(\mathcal G)$ denote the $M$ most recent ordered
elements of $\mathcal G$, with
$\operatorname{Tail}_{+\infty}(\mathcal G)=\mathcal G$. Let $M_s$ be the
minimum of the active date cap, the
window-derived cap $\lfloor W_{\max}/U\rfloor$, and, in default online mode,
the number of dates in the fixed-mode grid for the same $s$; an absent cap is
$+\infty$, and \code{full_online_history} omits the last term. The code returns
\[
 \mathcal R_s=\operatorname{Tail}_{M_s}(\mathcal G_s).
\]
Thus $A=1$ implies $(s-r)\bmod P=0$ for every
$r\in\mathcal R_s$; $A=0$ imposes no phase constraint.
In fixed mode, both $(r-L,r]$ and $(r,r+H]$ lie in $\Tzero$. In online mode,
\[
 r\geq L-1,\qquad r+H\leq s,
\]
so the neighbour future is already observed. It may overlap the query
lookback, but never the query target.

Without an explicit override, the code first counts the dates that a fixed
$\Tzero$ store would provide under the same value of $A$. The online
store then keeps that many \emph{most recent} eligible dates, subject to the
30,000-window cap. Thus the store slides forward without growing merely
because the query is later. \code{full_online_history} removes this implicit
date-count cap; explicit date and window caps still apply. Defaults are
$P=96$ samples for ETTm1, ETTm2, ETT\_T\_15T, and ETT\_L\_15T, and $P=24$
otherwise. Publication profiles disable period alignment so their datastore
origins rotate through seasonal phases. If period alignment is explicitly
enabled, the datastore stride must be a multiple of $P$.

\subsection{Fitting and leakage controls}

\begin{table}[H]
\centering
\small
\caption{What each chronological region may influence.}
\label{tab:fit-protocols}
\begin{tabularx}{\linewidth}{@{}lXXX@{}}
\toprule
Family & $\Tone$ & $\Ttwo$ & $\Tthree$ \\
\midrule
Convex mixtures & diagnostic fit & final fit on pooled $\Tpool$ & score once \\
Ridge baselines & fit all $\alpha$ & select $\alpha$; refit pooled & score once \\
Fixed-candidate CatBoost & fit up to 300 trees & select tree count; fresh pooled refit & score once \\
TS-IFA joint & optimize branches and one rooter & select and restore checkpoint & score once \\
TS-IFA meta & chronological support/query outer updates & fit active frozen-branch rooter & score once \\
Oracle diagnostics & -- & -- & target-aware appendix bound only \\
\bottomrule
\end{tabularx}
\end{table}

Ridge sufficient statistics are accumulated exactly in float64 chunks. Optional
sample caps affect fitting only; final $\Tthree$ scoring always uses every saved
sample. A future gate over a trainable candidate must train that candidate on
$\Tone$ and the gate on $\Ttwo$ without pooled refitting, unless out-of-fold
candidate predictions are introduced.

Dataset loaders discover \code{config.json} beside the selected CSV. Portable
fields live at the JSON top level and adaptation-only overrides under
\code{adaptation}. Project-scoped values override shared values, explicit run
arguments override both, and \code{drop_users} is additive. The resolved path
and applied keys are recorded in the extraction manifest; no file content hash
is created.

\section{TS-IFA}
\label{sec:tsifa}

TS-IFA is a lightweight neural adapter around the frozen backbone
$F:=F_\theta$. We again suppress $(s,u)$, so $X=X_{s,u}$ and $Y=Y_{s,u}$;
all following values are on the query scale:
\[
 p=F(X),\qquad p^c=F(X;\{X_j,Y_j\}_{j=1}^K),\qquad
 n_j=F(X_j),\qquad e_j=Y_j-n_j.
\]

\subsection{Candidate branches}

Define the learned maps
\begin{align*}
 A_r&:\reals^{L+H}\times(\reals^{L+H})^K\times(\reals^H)^K\to\reals^H,
 &M_r&:\reals^{2H}\to\reals^H,\\
 A_m&:\reals^L\times(\reals^L)^K\times(\reals^H)^K\to\reals^H,
 &M_m&:\reals^{2H}\to\reals^H.
\end{align*}
$A_r,A_m$ are multi-head cross-attention maps and $M_r,M_m$ are correction
MLPs. Residual attention compares query and neighbour \emph{states}:
\begin{align}
 q_r&=[X,p],& k_j^r&=[X_j,n_j],& v_j^r&=e_j,\\
 z_r&=A_r(q_r,\{k_j^r\},\{v_j^r\}),&
 c_r&=p+M_r([p,z_r]).
\end{align}
Memory attention reads realised neighbour horizons directly:
\[
 z_m=A_m(X,\{X_j\},\{Y_j\}),
 \qquad c_m=p+M_m([p,z_m]).
\]
Vanilla is always present, and any non-empty subset of
$\{p^c,c_r,c_m\}$ can be selected. For active candidate set $I$, the
unconstrained rooters operate on the non-vanilla deltas $d_j=c_j-p$, $j\in I$.
Fixed and learned transformed covariates are a
separate experiment rather than TS-IFA candidates.

\subsection{Neural and ridge rooters}

The neural rooter receives no handcrafted distances or dispersion features.
Its learned attention map is
$A_g:\reals^L\times(\reals^L)^K\times(\reals^L)^K\to\reals^{d_g}$,
where $d_g$ is the configured representation width:
\[
 z_g=A_g(X,\{X_j\},\{X_j\})\in\reals^{d_g}.
\]
For the learned type embedding $\tau_j\in\reals^{d_g}$ of active delta $d_j$,
let $\operatorname{LN}$ denote dimension-preserving layer normalization and
let $S_\omega:\reals^{2d_g+2H}\to\reals^R$ be the shared MLP scorer, where
$R=1$ for shared routing and $R=H$ for horizon routing. It returns
\begin{align}
 a_j&=S_\omega\!\left([
 \operatorname{LN}(z_g),\operatorname{LN}(p),
 \operatorname{LN}(d_j),\operatorname{LN}(\tau_j)
 ]\right)\in\reals^R,\\
 \widehat Y_{\rm neural}&=p+\sum_{j\in I}a_j\odot d_j.
\end{align}
Here $\omega$ denotes all neural-rooter parameters: $A_g$, its layer
normalizations, the type embeddings, and the shared scorer.
Unconstrained routing uses these signed delta weights. Optional softmax routing
prepends a vanilla logit, then normalizes vanilla and every active candidate to
non-negative weights summing to one. The shared form broadcasts one value over
all horizons; the horizon form retains one value per horizon index.

The ridge rooter uses the same active candidates. In the joint variant its free
routing matrix is $\nu\in\reals^{|I|\times R}$ and is updated
by gradient descent with the branches. Only the meta variant uses the following
closed-form solve.
For any non-empty fit set $\mathcal A$, let
$d_{i,j,h}$ be the horizon-$h$ delta of candidate $j$ on window $i$. Let
$n_{\mathcal A}=|\mathcal A|\geq1$. The differentiable meta-solve uses
$\epsilon_{\rm meta}=10^{-8}$ and the final fit uses
$\epsilon_{\rm fit}=10^{-12}$:
\[
 s_{j,h}^{(d,\mathcal A)}=
 \begin{cases}
  \sqrt{n_{\mathcal A}^{-1}\sum_{i\in\mathcal A}d_{i,j,h}^2
   +\epsilon_{\rm meta}^2},&\text{differentiable meta-solve},\\
  \max\!\left(\sqrt{n_{\mathcal A}^{-1}\sum_{i\in\mathcal A}d_{i,j,h}^2},
   \epsilon_{\rm fit}\right),&\text{final fit},
 \end{cases}
 \qquad \widetilde d_{i,j,h}=d_{i,j,h}/s_{j,h}^{(d,\mathcal A)}.
\]
For ridge penalty $\alpha>0$, the standardized coefficient matrix is
\[
 \Gamma^*_{\mathcal A}(\theta_b)=\arg\min_{\Gamma\in\reals^{3\times H}}
 \sum_{i\in\mathcal A}\sum_{h=1}^H
 \left(Y_{i,h}-p_{i,h}
  -\sum_{j=1}^{3}\Gamma_{j,h}\widetilde d_{i,j,h}\right)^2
 +\alpha\|\Gamma\|_F^2,
 \qquad \nu^*_{\mathcal A,j,h}(\theta_b)
 =\Gamma^*_{\mathcal A,j,h}(\theta_b)/s_{j,h}^{(d,\mathcal A)},
\]
where $\|\cdot\|_F$ is the Frobenius norm. This exact solve applies to
unconstrained meta-ridge; softmax meta-ridge uses differentiable gradient inner
updates. The unconstrained ridge prediction is
$p_{i,h}+\sum_{j=1}^{3}\nu_{j,h}d_{i,j,h}$. Both ridge variants have exactly
$|I|R$ routing parameters; $\alpha=10^{-2}$ applies to exact meta-ridge solves only.

\subsection{Training and initialization}

Let $\theta_b$ collect the parameters of $c_r$ and $c_m$, and let
$\mathcal B=\{c_r,c_m\}$. For minibatch index set $\mathcal A$, active-rooter
parameters $\psi\in\{\nu,\omega\}$, and coefficients $a_{i,j,h}$, define
\begin{align}
 \mathcal L_B^{\mathcal A}(\theta_b)
 &=\sum_{c\in\mathcal B}\nMSE_{\mathcal A}(c,Y)
 +\lambda_V\sum_{c\in\mathcal B}\nMSE_{\mathcal A}(c,p),\\
 \mathcal L_R^{\mathcal A}(\theta_b,\psi)
 &=\nMSE_{\mathcal A}(\widehat Y,Y)+\lambda_V\nMSE_{\mathcal A}(\widehat Y,p)
 +\frac{\lambda_2}{3|\mathcal A|H}\sum_{i,j,h}a_{i,j,h}^2
 +\frac{\lambda_S}{3|\mathcal A|(H-1)}\sum_{i,j,h<H}
 (a_{i,j,h+1}-a_{i,j,h})^2.
\end{align}
The last term is zero for $H=1$. Active values are
$\lambda_V=\lambda_2=\lambda_S=10^{-2}$ and $\lambda_B=0.1$.

The two joint variants use all of $\Tone$:
\begin{align}
 (\theta_b^*,\nu^*)&=\arg\min_{\theta_b,\nu}
 \mathcal L_R^{\Tone}(\theta_b,\nu)+\lambda_B\mathcal L_B^{\Tone}(\theta_b),
 &&\text{joint ridge},\\
 (\theta_b^*,\omega^*)&=\arg\min_{\theta_b,\omega}
 \mathcal L_R^{\Tone}(\theta_b,\omega)+\lambda_B\mathcal L_B^{\Tone}(\theta_b),
 &&\text{joint neural}.
\end{align}
AdamW updates both parameter blocks. Joint ridge never uses a closed-form solve.
The minimum-$\nMSE$ complete checkpoint on $\Ttwo$ is restored, then scored once
on $\Tthree$; no joint rooter is refitted on $\Ttwo$.

For meta learning, split $\Tone$ by whole query dates into earlier support
$\mathcal S_1$ and later query $\mathcal Q_1$. Meta ridge computes
\begin{align}
 \nu_e^*(\theta_b)&=\nu^*_{\mathcal S_1^{(e)}}(\theta_b),\\
 \theta_b^*&=\arg\min_{\theta_b}
 \mathcal L_R^{\mathcal Q_1^{(e)}}(\theta_b,\nu_e^*(\theta_b))
 +\lambda_B\mathcal L_B^{\mathcal Q_1^{(e)}}(\theta_b),
\end{align}
differentiating through each exact standardized $|I|\times|I|$ solve. It freezes
$\theta_b^*$ and fits $\nu^*_{\Ttwo}(\theta_b^*)$ in closed form.
Meta neural starts from learned $\omega^{(0)}=\omega$ and takes
\[
 \omega^{(k+1)}=\omega^{(k)}-\eta_{\rm in}
 \nabla_{\omega^{(k)}}\mathcal L_R^{\mathcal S_1^{(e)}}
 (\theta_b,\omega^{(k)}).
\]
Its query loss updates $\theta_b$ and the initialization $\omega$; first-order
MAML is the default. It then freezes $\theta_b$ and fits only $\omega$ on all
$\Ttwo$. Both meta variants score $\Tthree$ once. The default query fraction is
0.2, $K_{\rm in}=1$, and $\eta_{\rm in}=10^{-3}$.

Output-zero initialization of $M_r$, $M_m$, the neural scorer, and the free
joint-ridge matrix gives $c_r=c_m=p$, $a_j=0$, and $\widehat Y=p$ at step zero.
All AdamW stages use learning rate $10^{-5}$, weight decay $10^{-4}$, batch size
256, and 20,000 one-random-batch epochs outside smoke mode. The exact variant,
split, optimizer, regularizers, widths, caps, and seed enter the completion
signature.

\iffalse
Let $\theta_b$ collect the parameters of $c_r$ and $c_m$, and let
$\mathcal B=\{c_r,c_m\}$. Split $\Tone$ by whole query dates into an earlier
support set $\mathcal S_1$ and a later query set $\mathcal Q_1$; the active
query fraction is 0.2. Episode $e$ samples minibatches
$\mathcal S_1^{(e)}\subseteq\mathcal S_1$ and
$\mathcal Q_1^{(e)}\subseteq\mathcal Q_1$. Solve
$\nu^*_{\mathcal S_1^{(e)}}(\theta_b)$ exactly by the preceding active-candidate
horizon-wise systems. Starting from the learned neural initialization
$\omega^{(0)}=\omega$, take $K_{\rm in}$ support steps
\[
 \omega^{(k+1)}=\omega^{(k)}-\eta_{\rm in}
 \nabla_{\omega^{(k)}}\mathcal L_R^{\mathcal S_1^{(e)}}
 (\theta_b,\omega^{(k)}),\qquad k=0,\ldots,K_{\rm in}-1.
\]
The outer objective is evaluated only on the later query minibatch:
\begin{align}
 \mathcal L_{\rm meta}
 &=\lambda_{\rm ridge}\,
 \mathcal L_{\rm ridge}^{\mathcal Q_1^{(e)}}
 (\theta_b,\nu^*_{\mathcal S_1^{(e)}}(\theta_b))
 +\lambda_{\rm neural}\,
 \mathcal L_R^{\mathcal Q_1^{(e)}}(\theta_b,\omega^{(K_{\rm in})})
 +\lambda_B\mathcal L_B^{\mathcal Q_1^{(e)}}.
\end{align}
AdamW updates $\theta_b$ and the initialization $\omega$ from this objective.
The ridge term is differentiated through the exact solve. Neural adaptation uses
first-order MAML by default; full higher-order differentiation is configurable.

After meta-training, $\theta_b$ is frozen. Let $\mathcal Q_2$ denote the index
set of all $\Ttwo$ windows. The final ridge coefficients are
$\nu^*_{\mathcal Q_2}(\theta_b)$, while the neural rooter starts
from the learned $\omega$ and is fine-tuned on all $\Ttwo$. $\Tthree$ is scored
once with no selection or refit:
\[
 \Tone:(\mathcal S_1\rightarrow\mathcal Q_1)_{\rm bilevel}
 \longrightarrow \Ttwo:\{\nu^*,\omega_{\rm fit}\}
 \longrightarrow \Tthree:\text{evaluation}.
\]

Let $\widehat Y=\widehat Y_{\rm neural}$ and set
$\lambda_V=\lambda_2=\lambda_S=10^{-2}$. Every nMSE below suppresses the
conditioning ``$\mid X$'' from \cref{sec:forecasting} and averages that
per-window expression over the minibatch. Let $\mathcal S_{\rm mb}$ be the
minibatch index set, $B_{\rm mb}=|\mathcal S_{\rm mb}|$, and
$a_{i,j,h}$ the neural-rooter coefficient for its window $i$. The branch and
rooter losses are
\begin{align}
 \mathcal L_B
 &=\sum_{c\in\mathcal B}\nMSE(c,Y)
 +\lambda_V\sum_{c\in\mathcal B}\nMSE(c,p),\\
 \mathcal L_R
 &=\nMSE(\widehat Y,Y)+\lambda_V\nMSE(\widehat Y,p)
 +\frac{\lambda_2}{3B_{\rm mb}H}
 \sum_{i\in\mathcal S_{\rm mb}}\sum_{j=1}^{3}\sum_{h=1}^H a_{i,j,h}^2
 +\frac{\lambda_S}{3B_{\rm mb}(H-1)}
 \sum_{i\in\mathcal S_{\rm mb}}\sum_{j=1}^{3}\sum_{h=1}^{H-1}
 (a_{i,j,h+1}-a_{i,j,h})^2.
\end{align}
The last term applies when $H>1$ and is defined as zero when $H=1$.
The ridge query loss has the first two terms of $\mathcal L_R$, with ridge
regularization applied in its support solve. Active meta weights are
$\lambda_{\rm ridge}=\lambda_{\rm neural}=1$ and $\lambda_B=0.1$;
$K_{\rm in}=1$ and $\eta_{\rm in}=10^{-3}$. Outer and final-rooter AdamW use
 learning rate $10^{-5}$, weight decay $10^{-4}$, batch size 256, and 20,000
one-random-batch epochs outside smoke mode.
Output-zero initialization of $M_r$, $M_m$, and the shared rooter scorer gives
$c_r=c_m=p$, $a_j=0$, and $\widehat Y=p$ at step zero.
The result manifest stores the exact launcher training signature; a changed
split, optimizer, regularizer, model width, sample cap, or seed is incomplete
under the old signature and must be rerun.
\fi

\subsection{Parameter counts}

The full horizon-routing reference configuration uses four attention heads,
per-head dimension 32, 128-unit hidden layers, and three rooter deltas. Shared
routing or a branch subset reduces the rooter and branch counts. $K$ affects
compute and memory, not parameter count.

\begin{table}[H]
\centering
\small
\caption{Full horizon-routing TS-IFA reference parameter counts, reproduced by
the implementation and synchronized \code{config.json} artifacts.}
\label{tab:tsifa-counts}
\begin{tabular}{@{}rrrrrr@{}}
\toprule
$L$ & $H$ & Branches & Neural rooter & Full neural & Branches + ridge \\
\midrule
168 & 24  & 522,080 & 323,576 & 845,656 & 522,152 \\
336 & 48  & 645,056 & 397,424 & 1,042,480 & 645,200 \\
504 & 168 & 915,872 & 508,616 & 1,424,488 & 916,376 \\
512 & 64  & 759,808 & 471,232 & 1,231,040 & 760,000 \\
\bottomrule
\end{tabular}
\end{table}

At $512{:}64$, the default ridge rooter has $3H=192$ coefficients.

\begin{table}[H]
\centering
\small
\caption{Adaptation-parameter comparison at the released Cross-RAG
$512{:}64$ setting. Frozen backbone parameters are excluded from the first
numeric column. External counts follow the released implementations
\citep{tsragcode,crossragcode}.}
\label{tab:parameter-comparison}
\begin{tabular}{@{}lrrr@{}}
\toprule
Method & Trainable adaptation params. & $\times$ TS-IFA & Share of backbone \\
\midrule
TS-IFA neural & 1,231,040 & 1.00 & 1.03\% of Chronos-2 \\
TS-IFA branches + ridge & 760,000 & 0.62 & 0.63\% of Chronos-2 \\
TS-RAG added modules & 4,775,425 & 3.88 & 2.33\% of Bolt-base \\
Cross-RAG added modules & 8,712,192 & 7.08 & 4.25\% of Bolt-base \\
Full horizon ridge baseline, $K=3$ & 512 & 0.00042 & 0.00025\% of Bolt-base \\
One shared CatBoost gate & $\leq6,000$ proxy & $\leq0.0049$ & $\leq0.005$\% of Chronos-2 \\
\bottomrule
\end{tabular}
\end{table}

\section{Experimental design}
\label{sec:experiments}

\subsection{Datasets and model profiles}

\begin{table}[H]
\centering
\small
\caption{Dataset panels and their roles. All non-date variates are retained
for ETT, Weather, and Exchange.}
\label{tab:datasets}
\begin{tabularx}{\linewidth}{@{}p{.19\linewidth}p{.21\linewidth}X@{}}
\toprule
Role & Datasets & Preparation \\
\midrule
Primary screen
& Electricity, Traffic, Solar, exchange\_rate
& Hourly Electricity/Traffic; Solar is summed from 10-minute to hourly;
Exchange is daily. Electricity drops the nine configured channels. \\
Full homogeneous panels
& ETT\_T\_1H, ETT\_L\_1H, ETT\_T\_15T, ETT\_L\_15T
& Temperature and load are separated. Hourly panels use $P=24$; 15-minute
panels use $P=96$. \\
Mixed-quantity ablation
& ETTh1, ETTh2, ETTm1, ETTm2, Weather
& Original seven-variable ETT panels; Weather is hourly mean aggregation.
Cross-variable retrieval is reported only as an ablation. \\
\bottomrule
\end{tabularx}
\end{table}

\begin{table}[H]
\centering
\small
\caption{Resolved experiment profiles. A selected winner name fixes family,
formula, retrieval space, metric, $K$, and online mode.}
\label{tab:profiles}
\begin{tabularx}{\linewidth}{@{}p{.105\linewidth}p{.20\linewidth}p{.18\linewidth}p{.15\linewidth}X@{}}
\toprule
Profile & Datasets & $(L,H)$ & Backbone & Retrieval / methods \\
\midrule
\code{test} & Electricity & $504{:}168$ & Chronos-2 & raw Euclidean, $K=3$; smoke-only \\
\code{screen} & four primary & $168{:}24,336{:}48,504{:}168$ & Chronos-2
& raw/instance Euclidean, $K\in\{1,3\}$; 10 baselines and 8 gate/reference artifacts \\
\code{full} & primary + four separated ETT & default three & Chronos-2
& selected complete screen pipelines only \\
\code{ultra} & same as full & default three & Chronos-2, TabPFN-TS
& same winners and output root as full \\
\code{crossrag} & four primary & $512{:}64$ & Chronos-2, Chronos-Bolt
& winner at selected $K$ on Chronos-2; winner and released Cross-RAG at $K=15$ on Chronos-Bolt \\
\bottomrule
\end{tabularx}
\end{table}

Let
\[
 \mathcal D_P=\{\text{Electricity, Traffic, Solar, exchange\_rate}\},\quad
 \mathcal S=\{168{:}24,336{:}48,504{:}168\},
\]
\[
 \mathcal R=\{\text{raw, instance}\}\times\{\text{Euclidean}\}
 \times\{K=1,3\}\times\{\text{online}\},
\]
and let $\mathcal D_B$ be the seven stems in
\cref{tab:baseline-designs}. The $12=|\mathcal D_P||\mathcal S|$ screen
settings test, for every $r\in\mathcal R$,
\[
 \mathcal B_{\rm screen}
 =\{\code{cov_forecast},\code{avgy},\code{y_mean}\}
 \cup\{\code{<d>_ridge_shared}:d\in\mathcal D_B\},
\]
\[
 \mathcal G_{\rm screen}
 =\bigcup_{A\in\{\code{cov},\code{avgy}\}}
 \{\code{direct_<A>},\code{bayes_<A>_shared},
 \code{catboost_<A>_regressor_shared},\code{oracle_<A>_shared}\},
\]
where \code{direct_cov}=\code{cov_forecast} and
\code{direct_avgy}=\code{avgy}. Thus the baseline and gate screens contain
10 and 8 artifacts per retrieval pipeline; convex, delta, horizon, classifier,
and mixture forms are excluded.

Default publication datastore, pooled-adaptation-query, and evaluation-query
strides are $25/25/127$, respectively, with period alignment disabled. All
three strides are coprime with the 24-step hourly and 96-step quarter-hourly
periods, so sampled origins rotate through every seasonal phase. The datastore
cap is 30,000 windows. The test profile uses aligned strides $168/256/256$ and
2,048 store windows. Seed 1 is used throughout.

\subsection{Benchmarks and ablations}

\begin{table}[H]
\centering
\small
\caption{Selected baseline controls and one-family transformations. Retrieval
metric is Euclidean and mode is online in every row.}
\label{tab:selected-baselines}
\begin{tabularx}{\linewidth}{@{}p{.14\linewidth}p{.13\linewidth}p{.21\linewidth}p{.21\linewidth}X@{}}
\toprule
Stem & Space, $K$ & Horizon front & Convex front & Delta front \\
\midrule
\code{full} & instance, 3 & ridge shared / horizon & ridge / convex shared & ridge / delta shared \\
\code{y} & instance, 3 & ridge shared / horizon & ridge / convex shared & ridge / delta shared \\
\code{cov_y} & instance, 3 & ridge shared / horizon & ridge / convex shared & ridge / delta shared \\
\code{residual} & instance, 3 & ridge shared / horizon & ridge / convex shared & ridge / delta shared \\
\code{cov_avgy} & raw, 3 & ridge shared / horizon & ridge / convex shared & ridge / delta shared \\
\code{cov} & raw, 3 & ridge shared / horizon & ridge / convex shared & ridge / delta shared \\
\bottomrule
\end{tabularx}
\end{table}

Each row retains \code{<stem>_ridge_shared} as its control. The three fronts
replace only the suffix by \code{ridge_horizon}, \code{convex_shared}, or
\code{delta_ridge_shared}, respectively.

\begin{table}[H]
\centering
\small
\caption{Selected CatBoost controls and one-axis transformations.}
\label{tab:selected-catboost}
\begin{tabularx}{\linewidth}{@{}p{.12\linewidth}p{.17\linewidth}p{.23\linewidth}X@{}}
\toprule
$A$ & Space, $K$ & Control & Separate one-axis variants \\
\midrule
\code{avgy} & instance, 3 & regressor shared & classifier shared; regressor shared soft; regressor horizon \\
\code{cov} & instance, 1 & regressor shared & classifier shared; regressor shared soft; regressor horizon \\
\bottomrule
\end{tabularx}
\end{table}

For each $A$, the CatBoost front also includes matching direct, Bayes, and
oracle references. No transformation is crossed with another.

\begin{table}[H]
\centering
\small
\caption{Root submission fronts and tested model sets.}
\label{tab:slurm-workflow}
\begin{tabularx}{\linewidth}{@{}p{.31\linewidth}X@{}}
\toprule
Front & Model set / controlled axis \\
\midrule
\code{screen.slurm} & $\mathcal B_{\rm screen}$ and $\mathcal G_{\rm screen}$ on $\mathcal D_P\times\mathcal S\times\mathcal R$ \\
\code{benchmark.slurm} & complete selected screen pipelines; \code{full} uses Chronos-2, \code{ultra} adds TabPFN-TS \\
\code{k_ablation.slurm} & nine selected pipelines, $K\in\{1,3,5,10,15,20,100\}$ only \\
\code{h_ablation.slurm} & selected pipelines, $L=504$, $H\in\{24,168,504\}$ only \\
\code{l_ablation.slurm} & selected pipelines, $H=24$, $L\in\{24,168,504\}$ only \\
\code{horizon_baselines_ablation.slurm} & six controls in \cref{tab:selected-baselines}; shared versus horizon ridge \\
\code{convex_baselines_ablation.slurm} & same six; shared ridge versus shared convex \\
\code{delta_baselines_ablation.slurm} & same six; shared ridge versus shared delta-ridge \\
\code{catboost_ablation.slurm} & two controls in \cref{tab:selected-catboost}; three separate one-axis variants \\
\code{mixed_quantity_ablation.slurm} & selected pipelines transferred to the five mixed-quantity panels \\
\code{crossrag.slurm} & winner on Chronos-2 at its own $K$; winner and Cross-RAG on Chronos-Bolt at $(L,H,K)=(512,64,15)$ \\
\code{ts_ifa.slurm} & complete non-meta routing/branch grid over the selected test/full/ultra subset \\
\code{ts_ifa\_h\_ablation.slurm} & selected candidates at $L=504$ over three horizons \\
\code{ts_ifa\_l\_ablation.slurm} & selected candidates at $H=24$ over three lookbacks \\
\code{ts_ifa\_meta\_ridge.slurm} & meta-learning on textually selected ridge candidates \\
\code{ts_ifa\_meta\_neural.slurm} & meta-learning on textually selected neural candidates \\
\code{extraction.slurm} & standalone extraction smoke or TS-IFA-input front \\
\code{baselines.slurm} & standalone primary-baseline smoke front \\
\code{gates.slurm} & standalone primary-gate smoke front \\
\code{tables.slurm} & standalone result-table front \\
\bottomrule
\end{tabularx}
\end{table}

Earlier screen analysis selected six shared-ridge baselines, two hard shared
CatBoost regressors, and one shared Bayes gate; no convex or delta baseline was
eligible for that screen. Those historical synchronized artifacts lack enough
payload and configuration evidence for current-schema reuse and are archived.
They remain analysis evidence, but a current table or continued $K$ study must
rerun the affected extraction and result identities. Publication CatBoost fits
allow 300 trees; only the smoke profile uses two.

\subsection{Result identity, repetitions, and selection}
Extraction is a reusable workflow with identity
\begin{verbatim}
outputs/extraction/dataset/L_H/backbone/space/metric/k/mode/run_n/
\end{verbatim}
and vanilla uses \code{none/none/0/none}. Baselines, gates, Cross-RAG, and each
ablation use a separate family root with identity
\begin{verbatim}
outputs/adaptation/family/dataset/L_H/backbone/formula/space/metric/k/mode/run_n/
\end{verbatim}
TS-IFA instead orders \code{variant/routing_scope/routing_constraint/branch_set}
before \code{space/metric/k/mode}. Each config owns one directory. Optional row
and column config subsets affect only table layout.

Splits, budgets, optimizers, regularizers, fit caps, candidate lists, and
evaluation controls are pipeline configs in \code{manifest.json}; device and
scheduler placement are runtime-only. The recorded seed fixes all stochastic
choices. Repeated seeds, when used, are \code{seed_n} leaves. Mode is only a
test/full/ultra path-subset selector and is never a family, phase, path field,
or computation-signature field.

Every run uses \code{schema_version: 1}, declared configs, optional plain
provenance, launch attempts, required artifacts, and status
not-run/running/interrupted/completed. Run identity contains only the schema,
identity/model configs, pipeline/experiment configs, and seeds. It never
fingerprints or hashes code, Slurm files, data, weights, logs, outputs, checkpoints, or
directories; code/data changes are manual rerun decisions. The schema changes
only for a deliberate global contract break, and \code{new} creates another
repeat with unchanged parameters. The default \code{overwrite_exact} collision policy skips exact completion,
resumes exact interruption, and adds \code{run_n} for a changed pipeline.
Before a non-skipped attempt becomes running, schema-owned preparation removes
stale generated artifacts while preserving \code{manifest.json}, archived
manifest history, and completed \code{seed_n} leaves. Computation entry points
never delete the allocated \code{run_n} root. A manifest-less \code{run_n} is
an invalid partial directory and is reclaimed by the next allocation at that
index.
Finished work remains \code{ready}; completion is recorded only by the owning
Slurm workflow's final successful exit. Completed manifests are
authoritative and are not followed by file hashing or synchronization checks.
Tables support distinct/latest/average pipeline configurations and
selected/latest/distinct/average exact repeats. Explicit pipeline filters
select configurations and must match even with one run.
\code{SELECTED_RUNS.txt} records only automatic or pinned exact repeats, and
every report records requested filters and obtained manifests. Unmigratable legacy
material is isolated under
\code{archive/legacy_pre_schema_v1_2026-08-07} and is never read by
current tables.

After a successful root workflow, the launcher submits \code{publish.slurm}
with an \code{afterok} dependency. The handoff contains only the producer's exact
stdout/stderr and launch-tagged run/report roots. The publisher excludes
\code{*.pt}, \code{*.npy}, and \code{*.cbm}, commits only those paths on
\code{main}, sources \code{$HOME/codes/proxy.sh} using the shared two-line
\code{$HOME/codes/.secrets/proxy.credentials}, and pushes \code{origin/main} without pulling
or creating a pull request. Set \code{PUBLISH_RESULTS=false} to disable it; the
same script accepts \code{--job-id} for a manual retry.

Table construction accepts only current result manifests, keeps selected
pipelines family-qualified, and requires complete dataset--setting--model
coverage before writing any table. It also reads every selected current
\code{baseline_artifacts.pt}, expands the neighbour coefficients by rank, and
writes exact signed coefficient CSV files and imshow heatmaps; shared models
have one row and horizon models have one row per horizon. Direct baselines have
no fitted coefficients and are omitted. Complete extraction, baseline, and gate
configurations are reused by default, so resubmitting a profile or ablation
front rebuilds these tables and plots without refitting. For the $K$ ablation,
the same stage also overlays every candidate's average held-out nMSE improvement
against $K$ and marks the global optimum. The obsolete combined
\code{TS-IFA} directory is not a valid input; each current TS-IFA method owns
its isolated manifest and result folder.

For setting $c$, method $m$, method score $\nMSE_{m,c}$, and matching vanilla
nMSE $v_c$, the reported percentage improvement is
\[
 I_{m,c}=100\frac{v_c-\nMSE_{m,c}}{v_c}.
\]
The average table reports the unweighted mean of $I_{m,c}$ over complete
settings, not a percentage computed after pooling losses. If $n_c$ is the
number of held-out windows in setting $c$, window-level win rate is
\[
 W_{m,c}=100\frac1{n_c}\sum_{i=1}^{n_c}
 \mathbf1\!\left[
 \frac1H\sum_{h=1}^H(G_{i,h}-\widehat G_{m,i,h})^2
 <\frac1H\sum_{h=1}^H(G_{i,h}-V_{i,h})^2
 \right].
\]
Final claims use complete held-out $\Tthree$ scoring. Optimistic
\code{*_eval_fit} baselines and target-aware oracle gates are appendix
diagnostics only.

\section{Summary}

The project tests a hierarchy of increasing adaptation capacity around frozen
forecasting foundation models: direct analogue forecasts, ridge/convex/delta
linear families, conditional CatBoost gates, and TS-IFA. Every method
shares the same query dates, leakage-free online datastore, query-scale
neighbour signals, and held-out $\Tthree$ evaluation. This makes capacity,
retrieval choice, and validation protocol explicit rather than conflating them
with backbone quality.

\bibliographystyle{plainnat}
\bibliography{latex/adaptation_references}

\end{document}
